\section{対称性の使い方}\label{節:定積分の手法選択}

ここまでの方法は,どれも新しい原始関数を増やすものではない.定積分を,値の等しい別の積分に書き換え,2つを組にする方法である.では,与えられた定積分にどの置換を試せばよいのだろうか.最初に区間を見て,次にその区間を保つ変換を考えると,候補を絞れる.

\begin{center}
\begin{tabularx}{0.96\linewidth}{@{}lXX@{}}
\toprule
区間や式の特徴 & 試す変換 & 比べる部分\\
\midrule
$[-a,a]$ & $x\mapsto -x$ & 偶関数と奇関数の部分\\
$[a,b]$ & $x\mapsto a+b-x$ & 和が定数になるか,一次因子が消えるか\\
$\left[0,\frac{\pi}{2}\right]$ & $x\mapsto\frac{\pi}{2}-x$ & 正弦と余弦が交換されるか\\
長い区間と周期関数 & 整数周期だけ平行移動 & 完全な周期と残りの区間\\
$\left[\frac1b,b\right]$ と $\frac{\dd x}{x}$ & $x\mapsto\frac1x$ & 逆数にした関数との和\\
\bottomrule
\end{tabularx}
\end{center}

候補が見つかっても,対称性だけで最後まで解けるとは限らない.その場合は,対称性によって厄介な因子を1つ除き,残った部分を置換積分や部分積分に帰着させればよい.

\begin{bookpoint}
\textbf{POINT}\par
区間を保つ変換で $x$ などの厄介な因子を先に除き,残った積分を置換積分・部分積分・部分分数分解で計算する.
\end{bookpoint}

\begin{example}{反転と置換}
次の定積分を求めよ.
\[
I=\int_0^\pi\frac{x\sin x}{1+\sin^2x}\dd x.
\]
\end{example}

\paragraphbf{方針}
$\frac{\sin x}{1+\sin^2x}$ は $x$ を $\pi-x$ に変えても同じなので,反転で $x$ を $\frac{\pi}{2}$ に置き換える.残りは $u=\cos x$ と置き,部分分数分解する.

\paragraphbf{解答}
$I=\frac{\pi}{\sqrt2}\log(\sqrt2+1)$

\paragraphbf{解法}
$g(x)=\frac{\sin x}{1+\sin^2x}$ は $g(\pi-x)=g(x)$ を満たす.式\eqref{式:対称な重みの一次因子}より
\[
I=\frac{\pi}{2}J,\qquad
J=\int_0^\pi\frac{\sin x}{1+\sin^2x}\dd x
\]
である.\,$J$ では分子の $\sin x\dd x$ が $-\dd(\cos x)$ に対応する.\,$u=\cos x$ と置くと,\,$x=0,\pi$ に対して $u=1,-1$ であり,\,$1+\sin^2x=2-u^2$ だから
\[
J=\int_{-1}^1\frac{1}{2-u^2}\dd u
\]
となる.この分母は $(\sqrt2-u)(\sqrt2+u)$ に分かれるので,
\[
\frac{1}{2-u^2}
 =\frac{1}{2\sqrt2}\left(\frac{1}{\sqrt2-u}+\frac{1}{\sqrt2+u}\right)
\]
である.\,$-1\le u\le1$ では $\sqrt2-u$,\,$\sqrt2+u$ がともに正なので,
\[
\begin{aligned}
J&=\left[\frac{1}{2\sqrt2}
\log\frac{\sqrt2+u}{\sqrt2-u}\right]_{-1}^1\\
 &=\frac{1}{2\sqrt2}\left(\log\frac{\sqrt2+1}{\sqrt2-1}-\log\frac{\sqrt2-1}{\sqrt2+1}\right)\\
 &=\frac{1}{\sqrt2}\log\frac{\sqrt2+1}{\sqrt2-1}\\
 &=\frac{1}{\sqrt2}\log(\sqrt2+1)^2\\
 &=\sqrt2\log(\sqrt2+1)
\end{aligned}
\]
となる.4行目では $(\sqrt2+1)(\sqrt2-1)=1$ より $\frac{\sqrt2+1}{\sqrt2-1}=(\sqrt2+1)^2$ を使った.したがって
\[
\begin{aligned}
I&=\frac{\pi}{2}\cdot\sqrt2\log(\sqrt2+1)\\
 &=\frac{\pi}{\sqrt2}\log(\sqrt2+1)
\end{aligned}
\]
を得る.

\paragraphbf{部分積分から始めた場合}
もし最初に部分積分を使うなら,\,$\frac{\sin x}{1+\sin^2x}$ の原始関数が必要になり,さらにその対数を積分する項が現れる.反転によって $x$ を除いたため,必要な計算が1度の置換と部分分数分解で済んだのである.同じ式でも,不定積分では「微分の逆を作れるか」を見る.定積分ではそれに加えて,「区間を保つ置換で2つの式を補い合わせられるか」を見るとよい.
